\documentclass[oneside,senior,etd]{BYUPhys}

\usepackage{cmap} % Для корректной кодировки в pdf
\usepackage[utf8]{inputenc}
\usepackage{rotating}

\usepackage[russian]{babel}
\usepackage{amsfonts} % Пакеты для математических символов и теорем
\usepackage{amstext}
\usepackage{amssymb}
\usepackage{amsthm}
\usepackage{graphicx} % Пакеты для вставки графики
\usepackage{subfig}
\usepackage{color}
\usepackage[unicode]{hyperref}
\usepackage[nottoc]{tocbibind} % Для того, чтобы список литературы отображался в оглавлении
\usepackage{verbatim} % Для вставок заранее подготовленного текста в режиме as-is
\usepackage{listings}

\newcommand{\sectionbreak}{\clearpage} % Раздел с новой станицы

\usepackage{tikz}
\usepackage{pgfplots}
\usetikzlibrary{arrows, positioning, shapes}
\usepackage{adjustbox}

\usepackage{makecell}
\usepackage{booktabs}
\usepackage{boldline}

\usepackage{xcolor}
\usepackage{soul}
\usepackage{url}
\usepackage{multirow}
\usepackage{amsmath}

\usepackage{pifont}
\usepackage{indentfirst} % Делать отступ в начале первого параграфа

\usepackage{minted}
\usepackage[russian]{cleveref}

% \usepackage{showframe}
% \renewcommand\ShowFrameLinethickness{0.15pt}
% \renewcommand*\ShowFrameColor{\color{red}}

\newtheorem{theorem}{Theorem}[section]
\newtheorem{corollary}{Следствие}[theorem]
\newtheorem{lemma}[theorem]{Лемма}

% tikz styles
\tikzstyle{process} = [rectangle, minimum width=9em, minimum height=3em,
text centered, draw=black, align=center]
\tikzstyle{rounded} = [rectangle, minimum width=9em, minimum height=3em,
text centered, draw=black, align=center, rounded corners]
\tikzstyle{arrow} = [thick,->,>=stealth]
\tikzset{
  dot/.style = {circle, fill, minimum size=#1, inner sep=0pt, outer sep=0pt},
  dot/.default = 3pt % size of the circle diameter 
}

\renewcommand
\listingscaption{Листинг}

% Общие параметры листингов
\lstset{
  %frame=TB,
  showstringspaces=false,
  tabsize=4,
  basicstyle=\linespread{1.0}\tt\small, % делаем листинги компактнее
  breaklines=true,
  texcl=true, % русские буквы в комментах
  captionpos=b,
  aboveskip=\baselineskip,
  commentstyle=\tt
}

\newcommand{\todo}[1]{\textcolor{red}{#1}}

% DEBUG
% \usepackage{showframe}

\Faculty{Факультет вычислительной математики и кибернетики}
\Chair{Кафедра системного программирования}
% Лаборатория
% \Lab{~}
\Year{2025}
  \Month{Март}
  \City{Москва}
  \AuthorText{Автор:}
  \Author{Егоров Илья Георгиевич}
  \AuthorEng{Ilya Yegorov}
  \AcadGroup{527}

  \TitleTop{Формальные методы разработки программ}
  % Раскомментируйте, если нужны еще строчки названия
  % \TitleMiddle{}
  % \TitleBottom{Третья строка названия}
  % Uncomment if you need English title
  % \TitleTopEng{Thesis theme, first line}
  % \TitleMiddleEng{Thesis theme, second line}
  % \TitleBottomEng{Thesis theme, third line}

\docname{Практикум -- задание \textnumero 1}
  % \Advisor{}
  % \AdvisorDegree{\\}
  % Раскомментируйте, чтобы добавить научного консультанта
  % \Consultant{}
  % \ConsultantDegree{\\}

% Закомментируйте, если аннотация не нужна
% \Abstract{}
% Раскомментируйте, если нужна английская аннотация
% \AbstractEng{Abstract in English}

% Раскомментируйте, чтобы написать благодарности
% \Acknowledgments{Благодарности.}

%%%% DON'T change this. It is here because .sty does not support cyrillic cp properly %%%%
\TitlePageText{Титульная страница}
\University{Московский государственный университет имени М.В.Ломоносова}
\GrText{гр.}
\AdvisorText{Научный руководитель}
\ConsultantText{Научный консультант}
\AbstractText{Аннотация}
\AcknowledgmentsText{Благодарности}
\ListingText{Листинг}
\AlgorithmText{Алгоритм}

% Set PDF title and author
\hypersetup{
  pdftitle={\PDFTitle},
  pdfauthor={\PDFAuthor}
}

\begin{document}
\fixmargins
 \makepreliminarypages

\oneandhalfspace

\pdfbookmark[section]{\contentsname}{toc}
\tableofcontents

\section{Задание 1.1}

\subsection{Реализация, моделирование, требования}

Введём обозначения для максимального и минимального чисел в машинной арифметике:
\begin{align}
  INT\_MIN = -2^{31} \\
  INT\_MAX = 2^{31} - 1
\end{align}

Также введём несколько дополнительных предикатов для обозначения переполнения:
\begin{align}
  OF^+(x, y) \equiv x + y < INT\_MIN \vee x + y > INT\_MAX \\
  OF^-(x, y) \equiv x - y < INT\_MIN \vee x - y > INT\_MAX
\end{align}
\begin{lemma}\label{lem:of}
  При $INT\_MIN \leq x \leq INT\_MAX$ верны
\begin{align}
  OF^+(x, y) \equiv y < 0 \wedge x < INT\_MIN - y \vee y > 0 \wedge x > INT\_MAX - y \\
  OF^-(x, y) \equiv y < 0 \wedge x > INT\_MAX + y \vee y > 0 \wedge x < INT\_MIN + y
\end{align}
\end{lemma}

\begin{proof}
В силу
\begin{align}
  y \geq 0 \rightarrow x + y \geq x \geq INT\_MIN
\end{align}
Верно
\begin{align}
\begin{split}
  x + y < INT\_MIN &\equiv (y < 0 \vee y \geq 0) \wedge x + y < INT\_MIN \\
    &\equiv y < 0 \wedge x < INT\_MIN - y
\end{split}
\end{align}
Аналогично:
\begin{align}
  x + y < INT\_MIN &\equiv y > 0 \wedge x > INT\_MAX - y
\end{align}
Откуда и получаем доказываемую формулу. Аналогично для $OF^-(x, y)$
\end{proof}

Выпишем их отрицания:
\begin{align}
\begin{split}
  \overline{OF^+(x, y)} & \equiv INT\_MIN \leq x + y \leq INT\_MAX \\
    & \stackrel{(\ref{lem:of})}{\equiv}
    \overline{(y < 0 \wedge x < INT\_MIN - y)} \wedge \overline{(y > 0 \wedge x > INT\_MAX - y)} \\
    & \stackrel{(\ref{alt})}{\equiv}
    (y < 0 \rightarrow x \geq INT\_MIN - y) \wedge (y > 0 \rightarrow x \leq INT\_MAX - y)
\end{split} \\
\begin{split}
  \overline{OF^-(x, y)} & \equiv INT\_MIN \leq x - y \leq INT\_MAX \\
    & \stackrel{(\ref{lem:of})}{\equiv}
    \overline{(y < 0 \wedge x > INT\_MAX + y)} \wedge \overline{(y > 0 \wedge x < INT\_MIN + y)} \\
    & \stackrel{(\ref{alt})}{\equiv}
    (y < 0 \rightarrow x \leq INT\_MAX + y) \wedge (y > 0 \rightarrow x \geq INT\_MIN + y)
\end{split}
\end{align}

Дополнительно отметим, чтo указанные в \cref{lem:of} операции
никогда не приведут к целочисленному переполнению:
\begin{align}
\begin{split}
  INT\_MIN \leq x < 0   & \rightarrow   INT\_MIN < INT\_MIN - x \leq 0 \\
                        & \wedge        -1 \leq INT\_MAX + x < INT\_MAX
\end{split} \\
\begin{split}
  0 < x \leq INT\_MAX   & \rightarrow   INT\_MIN < INT\_MIN + x \leq -1 \\
                        & \wedge        0 \leq INT\_MAX - x < INT\_MAX
\end{split}
\end{align}

Везде далее доменами входных переменных $x = (x_1, x_2, x_3)$ и промежуточной
переменной $y$ являются целые
числа:
\begin{align}
  D_{x_1} = D_{x_2} = D_{x_3} = D_y = \mathbb{Z}, D_x = \mathbb{Z}^3
\end{align}


На \cref{p1src,p2src} представлен исходный код программ P1 и P2
соответственно, а на рисунках \cref{p1fc,p2fc}~--- их блок-схемы.

\begin{listing}[h]
\caption{Программа P1}
\label{p1src}
\begin{minted}[linenos=true]{c}
int comp(int x1, int x2, int x3) {
    int y = x1 - x3;
    return y + x2;
}
\end{minted}
\end{listing}

\begin{figure}[H]
\caption{Блок-схема программы P1}
\label{p1fc}
\centering
\begin{tikzpicture}[node distance=5em]
  % Nodes
  \node (start)     [rounded]                   {START: \\ $y \leftarrow 0$};
  \node (ofsub)     [rounded, below of=start]   {$OF^-(x_1, x_3)$};
  \node (assign)    [process, below of=ofsub]   {$y \leftarrow x_1 - x_3$};
  \node (ofadd)     [rounded, below of=assign]  {$OF^+(y, x_2)$};
  \node (halt)      [rounded, below of=ofadd]   {HALT: \\ $z \leftarrow y + x_2$};
  % Arrows
  \draw [arrow] (start) -- node[dot, pos=.7, label=0:A] {} coordinate (loopsub) (ofsub);
  \draw [arrow] (ofsub.west) -- +(-.5,0) |- node[left, pos=.25] {T} (loopsub-|ofsub.north);
  \draw [arrow] (ofsub) -- node[right] {F} (assign);
  \draw [arrow] (assign) -- node[dot, pos=.7, label=0:B] {} coordinate (loopadd) (ofadd);
  \draw [arrow] (ofadd.west) -| +(-.5,0) |- node[left, pos=.25] {T} (loopadd-|ofadd.north);
  \draw [arrow] (ofadd) -- node[right] {F} (halt);
\end{tikzpicture}
\end{figure}

\begin{listing}[H]
\caption{Программа P2}
\label{p2src}
\begin{minted}[linenos=true]{c}
#include <limits.h>

#define OFsub(x, y) (y < 0 && x > INT_MAX + y || y > 0 && x < INT_MIN + y)
#define OFadd(x, y) (y < 0 && x < INT_MIN - y || y > 0 && x > INT_MAX - y)
#define NOFsub(x, y) (!OFsub(x, y))
#define NOFadd(x, y) (!OFadd(x, y))

int comp(int x1, int x2, int x3) {
    int y;
    if (NOFsub(x1, x3)) {
        y = x1 - x3;
        return y + x2;
    }
    if (NOFadd(x1, x2)) {
        y = x1 + x2;
        return y - x3;
    }
    y = x2 - x3;
    return y + x1;
}
\end{minted}
\end{listing}

\begin{figure}[H]
\caption{Блок-схема программы P2}
\label{p2fc}
\centering
\begin{tikzpicture}[node distance=5em]
  % Nodes
  \node (start)     [rounded]                               {START: \\ $y \leftarrow 0$};
  % First try
  \node (ofsub1)    [rounded, below of=start]               {$OF^-(x_1, x_3)$};
  \node (assign1)   [process, right of=ofsub1, xshift=9em]  {$y \leftarrow x_1 - x_3$};
  \node (ofadd1)    [rounded, below of=assign1]             {$OF^+(y, x_2)$};
  \node (halt1)     [rounded, right of=ofadd1, xshift=9em, label=0:$(H_1)$]
                                                            {HALT: \\ $z \leftarrow y + x_2$};
  % Second try
  \node (ofadd2)    [rounded, below of=ofsub1, yshift=-5em] {$OF^+(x_1, x_2)$};
  \node (assign2)   [process, right of=ofadd2, xshift=9em]  {$y \leftarrow x_1 + x_2$};
  \node (ofsub2)    [rounded, below of=assign2]             {$OF^-(y, x_3)$};
  \node (halt2)     [rounded, right of=ofsub2, xshift=9em, label=0:$(H_2)$]
                                                            {HALT: \\ $z \leftarrow y - x_3$};
  % Third try
  \node (ofsub3)    [rounded, below of=ofadd2, yshift=-5em] {$OF^-(x_2, x_3)$};
  \node (assign3)   [process, right of=ofsub3, xshift=9em]  {$y \leftarrow x_2 - x_3$};
  \node (ofadd3)    [rounded, below of=assign3]             {$OF^+(y, x_1)$};
  \node (halt3)     [rounded, right of=ofadd3, xshift=9em, label=0:$(H_3)$]
                                                            {HALT: \\ $z \leftarrow y + x_1$};
  % Arrows
  % First try
  \draw [arrow] (start)     -- (ofsub1);
  \draw [arrow] (ofsub1)    -- node[left, pos=.25]  {T} node[dot, pos=.5,  label=0:B] {} (ofadd2);
  \draw [arrow] (ofsub1)    -- node[above]          {F} (assign1);
  \draw [arrow] (assign1)   -- node[dot, pos=.7, label=0:A] {} coordinate (loop) (ofadd1);
  \draw [arrow] (ofadd1.west) -| +(-.5,0) |- node[left, pos=.25] {T} (loop);
  \draw [arrow] (ofadd1)    -- node[above]          {F} (halt1);
  % Second try
  \draw [arrow] (ofadd2)    -- node[left, pos=.25]  {T} node[dot, pos=.75, label=0:D] {}
    coordinate (to3) (ofsub3);
  \draw [arrow] (ofadd2)    -- node[above]          {F} (assign2);
  \draw [arrow] (assign2)   -- node[dot, pos=.7, label=0:C] {} coordinate (loop) (ofsub2);
  \draw [arrow] (ofsub2.west) -| +(-.5,0) |- node[left, pos=.25] {T} (loop);
  \draw [arrow] (ofsub2)    -- node[above]          {F} (halt2);
  % Third try
  \draw [arrow] (ofsub3.west) -| +(-.5,0) |- node[left, pos=.25] {T} (to3);
  \draw [arrow] (ofsub3)    -- node[above]          {F} (assign3);
  \draw [arrow] (assign3)   -- node[dot, pos=.7, label=0:E] {} coordinate (loop) (ofadd3);
  \draw [arrow] (ofadd3.west) -| +(-.5,0) |- node[left, pos=.25] {T} (loop);
  \draw [arrow] (ofadd3)    -- node[above]          {F} (halt3);
\end{tikzpicture}
\end{figure}

Требования T1:
\begin{align}
\begin{split}
  \varphi_{T1}(x)   &\equiv INT\_MIN \leq x_1              \leq INT\_MAX \\
                    &\wedge INT\_MIN \leq x_2              \leq INT\_MAX \\
                    &\wedge INT\_MIN \leq x_3              \leq INT\_MAX \\
                    &\wedge INT\_MIN \leq x_1 - x_3        \leq INT\_MAX \\
                    &\wedge INT\_MIN \leq x_1 + x_2 - x_3  \leq INT\_MAX
\end{split} \\
  \psi(x)           &\equiv z = x_1 + x_2 - x_3
\end{align}

Требования T2:
\begin{align}
\begin{split}
  \varphi_{T2}(x)   &\equiv INT\_MIN \leq x_1              \leq INT\_MAX \\
                    &\wedge INT\_MIN \leq x_2              \leq INT\_MAX \\
                    &\wedge INT\_MIN \leq x_3              \leq INT\_MAX \\
                    &\wedge INT\_MIN \leq x_1 + x_2 - x_3  \leq INT\_MAX
\end{split} \\
  \psi(x)           &\equiv z = x_1 + x_2 - x_3
\end{align}

Докажем несколько вспомогательных утверждений.
\begin{lemma}\label{lem:sub}
\begin{align}
  \forall x \in D_x \ \varphi_{T1}(x) \rightarrow \overline{OF^-(x_1, x_3)}
\end{align}
\end{lemma}

\begin{proof}
\begin{align}
  & \varphi_{T1}(x) \rightarrow INT\_MIN \leq x_1 - x_3 \leq INT\_MAX \\
\begin{split}
  x_1 - x_3 \leq INT\_MAX \rightarrow x_1 \leq INT\_MAX + x_3
  \stackrel{(\ref{impl})}{\rightarrow} \\
  \rightarrow (x_3 < 0 \rightarrow x_1 \leq INT\_MAX + x_3)
\end{split} \\
\begin{split}
  x_1 - x_3 \geq INT\_MIN \rightarrow x_1 \geq INT\_MIN + x_3
  \stackrel{(\ref{impl})}{\rightarrow} \\
  \rightarrow (x_3 > 0 \rightarrow x_1 \geq INT\_MIN + x_3)
\end{split}
\end{align}
\end{proof}

\begin{lemma}\label{lem:full}
\begin{align}
  \forall x \in D_x \ \varphi_{T2}(x) \rightarrow \overline{OF^+(x_1 - x_3, x_2)} \\
  \forall x \in D_x \ \varphi_{T2}(x) \rightarrow \overline{OF^-(x_1 + x_2, x_3)} \\
  \forall x \in D_x \ \varphi_{T2}(x) \rightarrow \overline{OF^+(x_2 - x_3, x_1)}
\end{align}
\end{lemma}

\begin{proof}
  Аналогично \cref{lem:sub}.
\end{proof}

\begin{corollary}\label{cor:full}
\begin{align}
  \forall x \in D_x \ \overline{\varphi_{T2}(x) \wedge OF^+(x_1 - x_3, x_2)} \\
  \forall x \in D_x \ \overline{\varphi_{T2}(x) \wedge OF^-(x_1 + x_2, x_3)} \\
  \forall x \in D_x \ \overline{\varphi_{T2}(x) \wedge OF^+(x_2 - x_3, x_1)}
\end{align}
\end{corollary}

\begin{proof}
  Верно в силу \cref{alt}.
\end{proof}

\begin{lemma}\label{lem:possible}
\begin{align}
  \forall x \in D_x \ \varphi_{T2}(x) \rightarrow (OF^-(x_1, x_3)
  \rightarrow \overline{OF^+(x_1, x_2)} \vee \overline{OF^-(x_2, x_3)})
\end{align}
\end{lemma}

\begin{proof}
\begin{align}
  OF^-(x_1, x_3) &\equiv x_3 < 0 \wedge x_3 > INT\_MAX + x_3
    \vee x_1 > 0 \wedge x_1 < INT\_MIN + x_3 \\
  & \varphi_{T2}(x) \rightarrow x_1 + x_2 - x_3 \leq INT\_MAX \\
\begin{split}
    x_3 < 0 \wedge x_1 > INT\_MAX + x_3 \wedge x_1 + x_2 - x_3 \leq INT\_MAX \\
    \rightarrow x_1 + x_2 - x_3 \leq INT\_MAX < x_1 - x_3 \leq x_2 < 0
    \label{prf:minor}
\end{split} \\
  \stackrel{(\ref{prf:minor})}{\Rightarrow} \
  & INT\_MIN \leq x_1 + x_2 - x_3 < x_1 + x_2 < x_1 \leq INT\_MAX
  \label{prf:sub1} \\
  \stackrel{(\ref{prf:sub1})}{\Rightarrow} \
  & \overline{OF^+(x_1, x_2)} \label{prf:nof1} \\
  & \varphi_{T2}(x) \rightarrow x_1 + x_2 - x_3 \geq INT\_MIN \\
\begin{split}
    x_3 > 0 \wedge x_1 < INT\_MIN + x_3 \wedge x_1 + x_2 - x_3 \geq INT\_MIN \\
    \rightarrow x_1 - x_3 < INT\_MIN \leq x_1 + x_2 - x_3 \rightarrow x_2 > 0
    \label{prf:major}
\end{split} \\
  \stackrel{(\ref{prf:major})}{\Rightarrow} \
  & INT\_MIN < -INT\_MAX \leq -x3 < x_2 - x_3 < x_2 \leq INT\_MAX
  \label{prf:sub2} \\
  \stackrel{(\ref{prf:sub2})}{\Rightarrow} \
  & \overline{OF^-(x_2, x_3)} \label{prf:nof2} \\
  \stackrel{(\ref{prf:nof1}), (\ref{prf:nof2})}{\Rightarrow} \
  & \overline{OF^+(x_1, x_2)} \vee \overline{OF^-(x_2, x_3)}
\end{align}
\end{proof}

\begin{corollary}
\begin{align}
  \forall x \in D_x \ \varphi_{T2}(x) \rightarrow (\overline{OF^-(x_1, x_3)}
  \vee \overline{OF^+(x_1, x_2)} \vee \overline{OF^-(x_2, x_3)})
\end{align}
\end{corollary}

\subsection{P1-T1}

\subsubsection{Частичная корректность}

Положим:
\begin{align}
  \varphi(x) \equiv \varphi_{T1}(x)
\end{align}

Условия верификации:
\begin{align}
  SA: \ &\forall x \in D_x \ \varphi(x) \rightarrow p_A(x, r_{SA}(x, y)) \\
  ATA:\ &\forall x \in D_x \ \forall y \in D_y \ \varphi(x) \wedge p_A(x, y) \wedge R_{ATA}(x, y)
    \rightarrow p_A(x, r_{ATA}(x, y)) \\
  AFB:\ &\forall x \in D_x \ \forall y \in D_y \ \varphi(x) \wedge p_A(x, y) \wedge R_{AFB}(x, y)
    \rightarrow p_B(x, r_{AFB}(x, y)) \\
  BTB:\ &\forall x \in D_x \ \forall y \in D_y \ \varphi(x) \wedge p_B(x, y) \wedge R_{BTB}(x, y)
    \rightarrow p_B(x, r_{BTB}(x, y)) \\
  BFH:\ &\forall x \in D_x \ \forall y \in D_y \ \varphi(x) \wedge p_B(x, y) \wedge R_{BFH}(x, y)
    \rightarrow \psi(x, r_{BFH}(x, y))
\end{align}

Положим следующие индуктивные утверждения:
\begin{align}
  p_A(x, y) &\equiv \overline{OF^-(x_1, x_3)} \\
  p_B(x, y) &\equiv \overline{OF^+(y, x_2)} \wedge y = x_1 - x_3
\end{align}

Получаем:
\begin{align}
  SA: \ &\forall x \in D_x \ \varphi(x) \stackrel{(\ref{lem:sub})}{\rightarrow}
    \overline{OF^-(x_1,x_3)}\label{p1t1p:sa} \\
  ATA:\ &\forall x \in D_x \ \forall y \in D_y \ \varphi(x) \wedge \overline{OF^-(x_1, x_3)}
    \wedge OF^-(x_1, x_3) \rightarrow \overline{OF^-(x_1, x_3)}\label{p1t1p:ata} \\
\begin{split}
  AFB:\ &\forall x \in D_x \ \forall y \in D_y \ \varphi(x) \wedge \overline{OF^-(x_1, x_3)}
    \wedge \overline{OF^-(x_1, x_3)} \\
    &\qquad \stackrel{(\ref{lem:full})}{\rightarrow}
    \overline{OF^+(x_1 - x_3, x_2)} \wedge x_1 - x_3 = x_1 - x_3 \label{p1t1p:afb}
\end{split} \\
\begin{split}
  BTB:\ &\forall x \in D_x \ \forall y \in D_y \ \varphi(x) \wedge \overline{OF^+(y, x_2)}
    \wedge y = x_1 - x_3 \wedge OF^+(y, x_2) \\
    &\qquad \rightarrow \overline{OF^+(y, x_2)} \wedge y = x_1 - x_3 \label{p1t1p:bfb}
\end{split} \\
\begin{split}
  BFH:\ &\forall x \in D_x \ \forall y \in D_y \ \varphi(x) \wedge \overline{OF^+(y, x_2)} \wedge y = x_1 - x_3
        \wedge \overline{OF^+(y, x_2)} \\
        &\qquad \rightarrow \overline{OF^+(y, x_2)} \wedge x_1 - x_3 + x_2 = x_1 + x_2 - x_3
        \label{p1t1p:bfh}
\end{split}
\end{align}

Из \cref{p1t1p:sa,p1t1p:ata,p1t1p:afb,p1t1p:bfb,p1t1p:bfh} вытекает
частичная корректность:
\begin{align}\label{p1t1p}
  \{\varphi\}P1\{\psi\}
\end{align}

\subsubsection{Завершимость и полная корректность}\label{sec:p1t1t}

Пусть $W$~--- фундированное множество с частичным порядком $<$, а $u_A(x, y)$ и
$u_B(x, y)$~--- оценочные функции в точках A и B соответственно.

Условия корректности оценочных функций:
\begin{align}
  CORR_A:\ &\forall x \in D_x \ \forall y \in D_y \ \varphi(x) \wedge p_A(x, y)
    \rightarrow u_A(x, y) \in W \\
  CORR_B:\ &\forall x \in D_x \ \forall y \in D_y \ \varphi(x) \wedge p_B(x, y)
    \rightarrow u_B(x, y) \in W \\
\end{align}

Условия завершимости:
\begin{align}
  ATA: &\forall x \in D_x \forall y \in D_y \ \varphi(x) \wedge p_A(x, y) \wedge R_{ATA}(x, y)
        \rightarrow u_A(x, r_{ATA}(x, y)) < u_A(x, y) \\
  BTB: &\forall x \in D_x \forall y \in D_y \ \varphi(x) \wedge p_B(x, y) \wedge R_{BTB}(x, y)
        \rightarrow u_B(x, r_{BTB}(x, y)) < u_B(x, y)
\end{align}

Возьмём \footnote{На самом деле, можно взять произвольные $W$, $<$, $u_A(x, y)$,
$u_B(x, y)$}:
\begin{align}
  & W \equiv \mathbb{N}_0 \ \text{со стандартной операцией} \: < \\
  & u_A(x, y) = u_B(x, y) = 0
\end{align}

Получаем:
\begin{align}
  CORR_A:\ &\forall x \in D_x \ \forall y \in D_y \ \varphi(x) \wedge \overline{OF^-(x_1, x_3)}
    \rightarrow 0 \in \mathbb{N}_0 \label{p1t1t:corra} \\
  CORR_B:\ &\forall x \in D_x \ \forall y \in D_y \ \varphi(x) \wedge \overline{OF^+(y, x_2)}
    \wedge y = x_1 - x_3 \rightarrow 0 \in \mathbb{N}_0 \label{p1t1t:corrb} \\
  ATA:\ &\forall x \in D_x \ \forall y \in D_y \ \varphi(x) \wedge \overline{OF^-(x_1, x_3)}
    \wedge OF^-(x_1, x_3) \rightarrow 0 < 0 \label{p1t1t:ata} \\
  BTB:\ &\forall x \in D_x \ \forall y \in D_y \ \varphi(x) \wedge \overline{OF^+(y, x_2)}
    \wedge y = x_1 - x_3 \wedge OF^+(y, x_2) \rightarrow 0 < 0 \label{p1t1t:btb}
\end{align}

Из
\cref{p1t1p:sa,p1t1p:ata,p1t1p:afb,p1t1p:bfb,p1t1t:corra,p1t1t:corrb,p1t1t:ata,p1t1t:btb}
вытекает
завершимость:
\begin{align}\label{p1t1t}
  \{\varphi\}P1\{\mathbb{T}\}
\end{align}

Из \cref{p1t1p,p1t1t} вытекает полная корректность:
\begin{align}
  \langle\varphi\rangle P1 \langle\psi\rangle
\end{align}

\subsection{P1-T2}

\subsubsection{Частичная корректность}

Положим:
\begin{align}
  \varphi(x) \equiv \varphi_{T2}(x)
\end{align}

Положим следующие индуктивные утверждения:
\begin{align}
  p_A(x, y) &\equiv \mathbb{T} \\
  p_B(x, y) &\equiv y = x_1 - x_3
\end{align}

Получаем:
\begin{align}
  SA: \ &\forall x \in D_x \ \varphi(x) \rightarrow
        \mathbb{T} \label{p1t2p:sa} \\
  ATA:\ &\forall x \in D_x \ \forall y \in D_y \ \varphi(x) \wedge \mathbb{T}
    \wedge OF^-(x_1, x_3) \rightarrow \mathbb{T} \label{p1t2p:ata} \\
  AFB:\ &\forall x \in D_x \ \forall y \in D_y \ \varphi(x) \wedge \mathbb{T}
    \wedge \overline{OF^-(x_1, x_3)} \rightarrow x_1 - x_3 = x_1 - x_3 \label{p1t2p:afb} \\
  BTB:\ &\forall x \in D_x \ \forall y \in D_y \ \varphi(x) \wedge y = x_1 - x_3
    \wedge OF^+(y, x_2) \rightarrow y = x_1 - x_3 \label{p1t2p:bfb} \\
\begin{split}
  BFH:\ &\forall x \in D_x \ \forall y \in D_y \ \varphi(x) \wedge y = x_1 - x_3
    \wedge \overline{OF^+(y, x_2)} \\
    &\quad \rightarrow x_1 - x_3 + x_2 = x_1 + x_2 - x_3 \label{p1t2p:bfh}
\end{split}
\end{align}

Из \cref{p1t2p:sa,p1t2p:ata,p1t2p:afb,p1t2p:bfb,p1t2p:bfh} вытекает частичная
корректность:
\begin{align}
  \{\varphi\}P1\{\psi\}
\end{align}

\subsubsection{Завершимость и полная корректность}

Рассмотрим следующие входные параметры:
\begin{align}
  (x_1, x_2, x_3) = (INT\_MIN, 1, 1)
\end{align}

Тогда:
\begin{align}
  x_1 + x_2 - x_3 = INT\_MIN + 1 - 1 = INT\_MIN
\end{align}

Следовательно:
\begin{align}
  \varphi(x_1, x_2, x_3) \equiv \mathbb{T}
\end{align}

Но:
\begin{align}
\begin{split}
  OF^-(x_1, x_3)
    &\equiv 1 < 0 \wedge INT\_MIN > INT\_MAX + 1 \\
    &\vee 1 > 0 \wedge INT\_MIN < INT\_MIN + 1 \\
    &\equiv \mathbb{F} \wedge \mathbb{F} \vee \mathbb{T} \wedge \mathbb{T}
    \equiv \mathbb{F}
\end{split}
\end{align}

Т.е. программа зацикливается на $(x_1, x_2, x_3)$. Или, иначе говоря:
\begin{align}
  \exists x \in D_x: \varphi(x) \wedge M[P1](x) = \omega
\end{align}

Получаем:
\begin{align}
  \overline{\langle\varphi\rangle P1 \langle\psi\rangle}
\end{align}

\subsection{P2-T2}

\subsubsection{Частичная корректность}

Положим:
\begin{align}
  \varphi(x) \equiv \varphi_{T2}(x)
\end{align}

Условия верификации:
\begin{align}
  SFA:  \ &\forall x \in D_x \ \varphi(x) \wedge R_{SFA}(x, y)
    \rightarrow p_A(x, r_{SFA}(x, y)) \\
  STB:  \ &\forall x \in D_x \ \varphi(x) \wedge R_{STB}(x, y)
    \rightarrow p_B(x, r_{STB}(x, y)) \\
  AFH_1:\ &\forall x \in D_x \ \forall y \in D_y \ \varphi(x) \wedge p_A(x, y)
    \wedge R_{AFH_1}(x, y) \rightarrow \psi(x, r_{AFH_1}(x, y)) \\
  ATA:  \ &\forall x \in D_x \ \forall y \in D_y \ \varphi(x) \wedge p_A(x, y)
    \wedge R_{ATA}(x, y) \rightarrow p_A(x, r_{ATA}(x, y)) \\
  BFC:  \ &\forall x \in D_x \ \forall y \in D_y \ \varphi(x) \wedge p_B(x, y)
    \wedge R_{BFC}(x, y) \rightarrow p_C(x, r_{BFC}(x, y)) \\
  BTD:  \ &\forall x \in D_x \ \forall y \in D_y \ \varphi(x) \wedge p_B(x, y)
    \wedge R_{BTD}(x, y) \rightarrow p_D(x, r_{BTD}(x, y)) \\
  CFH_2:\ &\forall x \in D_x \ \forall y \in D_y \ \varphi(x) \wedge p_C(x, y)
    \wedge R_{CFH_2}(x, y) \rightarrow \psi(x, r_{CFH_2}(x, y)) \\
  CTC:  \ &\forall x \in D_x \ \forall y \in D_y \ \varphi(x) \wedge p_C(x, y)
    \wedge R_{CTC}(x, y) \rightarrow p_C(x, r_{CTC}(x, y)) \\
  DFE:  \ &\forall x \in D_x \ \forall y \in D_y \ \varphi(x) \wedge p_D(x, y)
    \wedge R_{DFE}(x, y) \rightarrow p_E(x, r_{DFE}(x, y)) \\
  DTD:  \ &\forall x \in D_x \ \forall y \in D_y \ \varphi(x) \wedge p_D(x, y)
    \wedge R_{DTD}(x, y) \rightarrow p_D(x, r_{DTD}(x, y)) \\
  EFH_3:\ &\forall x \in D_x \ \forall y \in D_y \ \varphi(x) \wedge p_E(x, y)
    \wedge R_{EFH_3}(x, y) \rightarrow \psi(x, r_{EFH_3}(x, y)) \\
  ETE:  \ &\forall x \in D_x \ \forall y \in D_y \ \varphi(x) \wedge p_E(x, y)
    \wedge R_{ETE}(x, y) \rightarrow p_E(x, r_{ETE}(x, y))
\end{align}

Положим следующие индуктивные утверждения:
\begin{align}
  p_A(x, y) &\equiv \overline{OF^+(y, x_2)} \wedge y = x_1 - x_3 \\
  p_B(x, y) &\equiv \overline{OF^+(x_1, x_2)} \vee \overline{OF^-(x_2, x_3)} \\
  p_C(x, y) &\equiv \overline{OF^-(y, x_3)} \wedge y = x_1 + x_2 \\
  p_D(x, y) &\equiv \overline{OF^-(x_2, x_3)} \\
  p_E(x, y) &\equiv \overline{OF^+(y, x_1)} \wedge y = x_2 - x_3
\end{align}

Получаем:
\begin{align}
  SFA:  \ &\forall x \in D_x \ \varphi(x) \wedge \overline{OF^-(x_1, x_3)}
    \stackrel{(\ref{lem:full})}{\rightarrow}
    \overline{OF^+(x_1 - x_3, x_2)} \wedge x_1 - x_3 = x_1 - x_3 \label{p2t2p:sfa} \\
  STB:  \ &\forall x \in D_x \ \varphi(x) \wedge OF^-(x_1, x_3)
    \stackrel{(\ref{lem:possible})}{\rightarrow}
    \overline{OF^+(x_1, x_2)} \vee \overline{OF^-(x_2, x_3)} \label{p2t2p:stb} \\
\begin{split}
  AFH_1:\ &\forall x \in D_x \ \forall y \in D_y \ \varphi(x) \wedge
    \overline{OF^+(y, x_2)} \wedge y = x_1 - x_3 \wedge \overline{OF^+(y, x_2)} \\
    &\quad \rightarrow x_1 - x_3 + x_2 = x_1 + x_2 - x_3 \label{p2t2p:afh1}
\end{split} \\
\begin{split}
  ATA:  \ &\forall x \in D_x \ \forall y \in D_y \ \varphi(x) \wedge
    \overline{OF^+(y, x_2)} \wedge y = x_1 - x_3 \wedge OF^+(y, x_2) \\
    &\quad \rightarrow \overline{OF^+(y, x_2)} \wedge y = x_1 - x_3 \label{p2t2p:ata}
\end{split} \\
\begin{split}
  BFC:  \ &\forall x \in D_x \ \forall y \in D_y \ \varphi(x) \wedge
    (\overline{OF^+(x_1, x_2)} \vee \overline{OF^-(x_2, x_3)}) \wedge
    \overline{OF^+(x_1, x_2)} \\
    &\quad \stackrel{(\ref{lem:full})}{\rightarrow}
    \overline{OF^-(x_1 + x_2, x_3)} \wedge x_1 + x_2 = x_1 + x_2 \label{p2t2p:bfc}
\end{split} \\
\begin{split}
  BTD:  \ &\forall x \in D_x \ \forall y \in D_y \ \varphi(x) \wedge
    (\overline{OF^+(x_1, x_2)} \vee \overline{OF^-(x_2, x_3)}) \wedge OF^+(x_1, x_2) \\
    &\quad \stackrel{(\ref{hack})}{\rightarrow}
    \overline{OF^-(x_2, x_3)} \label{p2t2p:btd}
\end{split} \\
\begin{split}
  CFH_2:\ &\forall x \in D_x \ \forall y \in D_y \ \varphi(x) \wedge
    \overline{OF^-(y, x_3)} \wedge y = x_1 + x_2 \wedge \overline{OF^-(y, x_3)} \\
    &\quad \rightarrow x_1 + x_2 - x_3 = x_1 + x_2 - x_3 \label{p2t2p:cfh2}
\end{split} \\
\begin{split}
  CTC:  \ &\forall x \in D_x \ \forall y \in D_y \ \varphi(x) \wedge
    \overline{OF^-(y, x_3)} \wedge y = x_1 + x_2 \wedge OF^-(y, x_3) \\
    &\quad \rightarrow \overline{OF^-(y, x_3)} \wedge y = x_1 + x_2 \label{p2t2p:ctc}
\end{split} \\
\begin{split}
  DFE:  \ &\forall x \in D_x \ \forall y \in D_y \ \varphi(x) \wedge
    \overline{OF^-(x_2, x_3)} \wedge \overline{OF^-(x_2, x_3)} \\
    &\quad \stackrel{(\ref{lem:full})}{\rightarrow}
    \overline{OF^+(x_2 - x_3, x_1)} \wedge x_2 - x_3 = x_2 - x_3 \label{p2t2p:dfe}
\end{split} \\
  DTD:  \ &\forall x \in D_x \ \forall y \in D_y \ \varphi(x) \wedge
    \overline{OF^-(x_2, x_3)} \wedge OF^-(x_2, x_3)
    \rightarrow \overline{OF^-(x_2, x_3)} \label{p2t2p:dtd} \\
\begin{split}
  EFH_3:\ &\forall x \in D_x \ \forall y \in D_y \ \varphi(x) \wedge
    \overline{OF^+(y, x_1)} \wedge y = x_2 - x_3 \wedge \overline{OF^+(y, x_1)} \\
    &\quad \rightarrow x_2 - x_3 + x_1 = x_1 + x_2 - x_3 \label{p2t2p:efh3}
\end{split}
\end{align}
\begin{align}
\begin{split}
  ETE:  \ &\forall x \in D_x \ \forall y \in D_y \ \varphi(x) \wedge
    \overline{OF^+(y, x_1)} \wedge y = x_2 - x_3 \wedge OF^+(y, x_1) \\
    &\quad \rightarrow \overline{OF^+(y, x_1)} \wedge y = x_2 - x_3 \label{p2t2p:ete} 
\end{split}
\end{align}

Из
\cref{p2t2p:sfa,p2t2p:stb,p2t2p:afh1,p2t2p:ata,p2t2p:bfc,p2t2p:btd,p2t2p:cfh2,p2t2p:ctc,p2t2p:dfe,p2t2p:dtd,p2t2p:efh3,p2t2p:ete}
вытекает частичная корректность:
\begin{align}\label{p2t2p}
  \{\varphi\}P2\{\psi\}
\end{align}

\subsubsection{Завершимость и полная корректность}

Пусть $W$~--- фундированное множество с частичным порядком $<$, а $u_A(x, y)$,
$u_C(x, y)$, $u_D(x, y)$ и $u_E(x, y)$ ~--- оценочные функции в точках A, C, D и
E соответственно.

Условие корректоности оценочной функции:
\begin{align}
  CORR_A:\ &\forall x \in D_x \ \forall y \in D_y \ \varphi(x) \wedge p_A(x, y)
            \rightarrow u_A(x, y) \in W \\
  CORR_C:\ &\forall x \in D_x \ \forall y \in D_y \ \varphi(x) \wedge p_C(x, y)
            \rightarrow u_C(x, y) \in W \\
  CORR_D:\ &\forall x \in D_x \ \forall y \in D_y \ \varphi(x) \wedge p_D(x, y)
            \rightarrow u_D(x, y) \in W \\
  CORR_E:\ &\forall x \in D_x \ \forall y \in D_y \ \varphi(x) \wedge p_E(x, y)
            \rightarrow u_E(x, y) \in W \\
\end{align}

Условие завершимости:
\begin{align}
\begin{split}
  ATA:\ &\forall x \in D_x \ \forall y \in D_y \ \varphi(x) \wedge p_C(x, y)
    \wedge R_{ATA}(x, y) \\
    &\quad \rightarrow u_A(x, r_{ATA}(x, y)) < u_A(x, y)
\end{split} \\
\begin{split}
  CTC:\ &\forall x \in D_x \ \forall y \in D_y \ \varphi(x) \wedge p_C(x, y)
    \wedge R_{CTC}(x, y) \\
    &\quad \rightarrow u_C(x, r_{CTC}(x, y)) < u_C(x, y)
\end{split} \\
\begin{split}
  DTD:\ &\forall x \in D_x \ \forall y \in D_y \ \varphi(x) \wedge p_C(x, y)
    \wedge R_{DTD}(x, y) \\
    &\quad \rightarrow u_D(x, r_{DTD}(x, y)) < u_D(x, y)
\end{split} \\
\begin{split}
  ETE:\ &\forall x \in D_x \ \forall y \in D_y \ \varphi(x) \wedge p_C(x, y)
    \wedge R_{ETE}(x, y) \\
    &\quad \rightarrow u_E(x, r_{ETE}(x, y)) < u_E(x, y)
\end{split}
\end{align}

Аналогично P1 (см. \cref{sec:p1t1t}) возьмём \footnote{Здесь также можно взять
произвольные $W$, $<$, $u_C(x, y)$}:
\begin{align}
  & W \equiv \mathbb{N}_0 \ \text{со стандартной операцией} \: < \\
  & u_A(x, y) = u_B(x, y) = u_C(x, y) = u_D(x, y) = 0
\end{align}

Получаем:
\begin{align}
  CORR_A:\ &\forall x \in D_x \ \forall y \in D_y \ \varphi(x) \wedge
            \overline{OF^+(y, x_2)} \wedge y = x_1 - x_3
            \rightarrow 0 \in \mathbb{N}_0 \label{p2t2t:ca} \\
  CORR_C:\ &\forall x \in D_x \ \forall y \in D_y \ \varphi(x) \wedge
            \overline{OF^-(y, x_3)} \wedge y = x_1 + x_2
            \rightarrow 0 \in \mathbb{N}_0 \label{p2t2t:cc} \\
  CORR_D:\ &\forall x \in D_x \ \forall y \in D_y \ \varphi(x) \wedge
            \overline{OF^-(x_2, x_3)}
            \rightarrow 0 \in \mathbb{N}_0 \label{p2t2t:cd} \\
  CORR_E:\ &\forall x \in D_x \ \forall y \in D_y \ \varphi(x) \wedge
            \overline{OF^+(y, x_1)} \wedge y = x_2 - x_3
            \rightarrow 0 \in \mathbb{N}_0 \label{p2t2t:ce} \\
\begin{split}
  ATA:\ &\forall x \in D_x \ \forall y \in D_y \ \varphi(x) \wedge
    \overline{OF^+(y, x_2)} \wedge y = x_1 - x_3 \wedge OF^+(y, x_2) \\
    &\quad \rightarrow 0 < 0 \label{p2t2t:ata}
\end{split} \\
\begin{split}
  CTC:\ &\forall x \in D_x \ \forall y \in D_y \ \varphi(x) \wedge
    \overline{OF^-(y, x_3)} \wedge y = x_1 + x_2 \wedge OF^-(y, x_3) \\
    &\quad \rightarrow 0 < 0 \label{p2t2t:ctc}
\end{split} \\
\begin{split}
  DTD:\ &\forall x \in D_x \ \forall y \in D_y \ \varphi(x) \wedge
    \overline{OF^-(x_2, x_3)} \wedge OF^-(x_2, x_3) \\
    &\quad \rightarrow 0 < 0 \label{p2t2t:dtd}
\end{split} \\
\begin{split}
  ETE:\ &\forall x \in D_x \ \forall y \in D_y \ \varphi(x) \wedge
    \overline{OF^+(y, x_1)} \wedge y = x_2 - x_3 \wedge OF^+(y, x_1) \\
    &\quad \rightarrow 0 < 0 \label{p2t2t:ete}
\end{split}
\end{align}

Из
\cref{p2t2t:ca,p2t2t:cc,p2t2t:cd,p2t2t:ce,p2t2t:ata,p2t2t:ctc,p2t2t:dtd,p2t2t:ete}
вытекает завершимость:
\begin{align}\label{p2t2t}
  \{\varphi\}P2\{\mathbb{T}\}
\end{align}

Из \cref{p2t2p,p2t2t} вытекает полная корректность:
\begin{align}\label{p2t2tt}
  \langle\varphi\rangle P2 \langle\psi\rangle
\end{align}


\subsection{P2-T1}

Заметим, что:
\begin{align}
  \varphi_{T1}(x) \rightarrow \varphi_{T2}(x)
\end{align}

Тогда:
\begin{align}
  \{\varphi_{T2}\} P2 \{\psi\} &\rightarrow
  \{\varphi_{T1}\} P2 \{\psi\} \label{pweak} \\
  \langle\varphi_{T2}\rangle P2 \langle\psi\rangle &\rightarrow
  \langle\varphi_{T1}\rangle P2 \langle\psi\rangle \label{tweak}
\end{align}

Из \cref{p2t2p,p2t2tt,pweak,tweak} имеем:
\begin{align}
  \{\varphi_{T1}\} &P2 \{\psi\} \\
  \langle\varphi_{T1}\rangle &P2 \langle\psi\rangle
\end{align}

\section{Задание 1.2}

Условия (см. блок-схему на \cref{st2fc}):
\begin{align}
  & V = \{x, y_1, y_2, y_3, z\} \nonumber \\
  & D_x = D_{y_1} = D_{y_2} = D_{y_3} = D_z = \mathbb{Z} \nonumber \\
  & y = (y_1, y_2, y_3), D_y = \mathbb{Z}^3 \nonumber \\
  & \varphi(x) \equiv x \geq 0 \\
  & \psi(x, z) \equiv z^3 \leq x < (z + 1)^3
\end{align}

\begin{figure}[H]
\caption{Блок-схема программы задания 1.2}
\label{st2fc}
\centering
\begin{tikzpicture}[node distance=5em]
  % Nodes
  \node (start) [rounded]                   {START: \\ $(y_1, y_2, y_3) \leftarrow (0, 0, 1)$};
  \node (y2ass) [process, below of=start]   {$y_2 \leftarrow y_2 + y_3$};
  \node (check) [rounded, below of=y2ass]   {$y_2 > x$};
  \node (y1ass) [process, below of=check, xshift=-6em]
                                            {$y_1 \leftarrow y_1 + 1$};
  \node (y3ass) [process, below of=y1ass]   {$y_3 \leftarrow y_3 + 6y_1$};
  \node (halt)  [rounded, below of=check, xshift=+6em]
                                            {HALT: \\ $z \leftarrow y_1$};
  % Arrows
  \draw [arrow] (start) -- coordinate (loop) (y2ass);
  \draw [arrow] (y2ass) -- node[dot, pos=.5, label=0:A] {} (check);
  \draw [arrow] (check.east) -| node[above, pos=.25] {T} (halt);
  \draw [arrow] (check.west) -| node[above, pos=.25] {F} (y1ass);
  \draw [arrow] (y1ass) -- (y3ass);
  \draw [arrow] (y3ass.south) -- +(0,-.5) -- +(-3,-.5) |- (loop);
\end{tikzpicture}
\end{figure}

Условия верификации:
\begin{align}
  SA: \ &\forall x \in D_x \ \varphi(x) \rightarrow p_A(x, r_{SA}(x, y)) \\
  AFA:\ &\forall x \in D_x \ \forall y \in D_y \ \varphi(x) \wedge p_A(x, y) \wedge R_{AFA}(x, y) \rightarrow p_A(x, r_{AFA}(x, y)) \\
  ATH:\ &\forall x \in D_x \ \forall y \in D_y \ \varphi(x) \wedge p_A(x, y) \wedge R_{ATH}(x, y) \rightarrow \psi(x, r_{ATH}(x, y))
\end{align}

Положим следующее индуктивное утверждение:
\begin{align}
  p_A(x, y) \equiv y_1^3 \leq x \wedge y_2 = (y_1 + 1)^3 \wedge y_3 = 3y_1^2 + 3y_1 + 1
\end{align}

%элементы $p_A(x, y_1 + 1, y_2 + y_3 + 6y_1 + 6, y_3 + 6y_1 + 1)$:

Получаем:
\begin{align}
  SA: \ &\forall x \in D_x \ x \geq 0 \rightarrow 0^3 \leq x \wedge 1^3 = (0 + 1)^3
  \wedge 1 = 0 + 0 + 1 \label{st2sa} \\
\begin{split}
  AFA:\ &\forall x \in D_x \ \forall y \in D_y \ x \geq 0 \wedge y_1^3 \leq x
    \wedge y_2 = (y_1 + 1)^3 \wedge y_3 = 3y_1^2 + 3y_1 + 1 \\
    &\quad \wedge y_2 \leq x
    \rightarrow (y_1 + 1)^3 \leq x \wedge y_2 + y_3 + 6y_1 + 6 = (y_1 + 2)^3 \\
    &\quad \wedge y_3 + 6y_1 + 6 = 3(y_1 + 1)^2 + 3(y_1 + 1) + 1 \label{st2afa}
\end{split} \\
\begin{split}
  ATH:\ &\forall x \in D_x \ \forall y \in D_y \ x \geq 0 \wedge y_1^3 \leq x
    \wedge y_2 = (y_1 + 1)^3 \wedge y_3 = 3y_1^2 + 3y_1 + 1 \\
    &\quad \wedge y_2 > x \rightarrow y_1^3 \leq x < (y_1 + 1)^3 \label{st2ath}
\end{split}
\end{align}

Покажем \cref{st2afa}:
\begin{align}
  & y_2 = (y_1 + 1)^3 \wedge y_2 \leq x \rightarrow (y_1 + 1)^3 \leq x
  \label{st2fa1proof} \\
\begin{split}
  & y_2 + y_3 + 6y_1 + 6 = y_1^3 + 3y_1^2 + 3y_1 + 1 + 3y_1^2 + 3y_1 + 1 + 6y_1 + 6 \\
  & \qquad = y_1^3 + 6y_1^2 + 12y_1 + 8 = (y_1 + 2)^3
\end{split} \\
\begin{split}
  & y_3 + 6y_1 + 6 = 3y_1^2 + 3y_1 + 1 + 6y_1 + 6 = 3y_1^2 + 9y_1 + 7 \\
  & \qquad = 3y_1^2 + 6y_1 + 3 + 3y_1 + 3 + 1 = 3(y_1 + 1)^2 + 3(y_1 + 1) + 1
\end{split}
\end{align}

Аналогично \cref{st2fa1proof} показывается \cref{st2ath}.

Из \cref{st2sa,st2afa,st2ath} вытекает частичная корректность:
\begin{align}
  \{\varphi\}P\{\psi\}
\end{align}

\section{Задание 1.3}

Условия (см. блок-схему на \cref{st3fc}):
\begin{align}
  & V = \{x, y_1, y_2, y_3, z\} \nonumber \\
  & D_x = D_{y_1} = D_{y_2} = D_{y_3} = D_z = \mathbb{Z} \nonumber \\
  & y = (y_1, y_2, y_3), D_y = \mathbb{Z}^3 \nonumber \\
  & \varphi(x) \equiv x > 1
\end{align}

\begin{figure}[H]
\caption{Блок-схема программы задания 1.3}
\label{st3fc}
\centering
\begin{tikzpicture}[node distance=5em]
  % Nodes
  \node (start)     [rounded]                   {START: \\ $(y_1, y_2, y_3)
                                                \leftarrow (x, 1, x)$};
  \node (checkles)  [rounded, below of=start]   {$y_2 < y_3$};
  \node (y1ass)     [process, below of=checkles, xshift=-6em]
                                                {$(y_1, y_2) \leftarrow
                                                (y_1 + x, y_2 + 1)$};
  \node (checkequ)  [rounded, below of=checkles, xshift=+11em]
                                                {$y_3 = x$};
  \node (y3ass)     [process, below of=checkequ, xshift=-6em]
                                                {$y_3 \leftarrow y_1$};
  \node (halt)      [rounded, below of=checkequ, xshift=+6em]
                                            {HALT: \\ $z \leftarrow y_1$};
  % Arrows
  \draw [arrow] (start) -- node[dot, pos=.7, label=0:A] {} coordinate (loop) (checkles);
  \draw [arrow] (checkles.west) -| node[above, pos=.25] {T} (y1ass);
  \draw [arrow] (checkles.east) -| node[above, pos=.25] {F} (checkequ);
  \draw [arrow] (checkequ.west) -| node[above, pos=.25] {T} (y3ass);
  \draw [arrow] (checkequ.east) -| node[above, pos=.25] {F} (halt);
  \draw [arrow] (y3ass.south) -- +(0,-.5) -- coordinate (bottom) +(-8,-.5) |- (loop);
  \draw [-]     (y1ass.south) -- (y1ass|-bottom);
\end{tikzpicture}
\end{figure}

Пусть $W$~--- фундированное множество с частичным порядком $<$, а $p_A(x, y)$ и
$u_A(x, y)$~--- индуктивное утверждение и оценочная функция в точкe A
соответственно.

Условия верификации:
\begin{align}
  SA:  \ &\forall x \in D_x \ \forall y \in D_y \ \varphi(x)
    \rightarrow p_A(x, r_{SA}(x, y)) \\
  ATA: \ &\forall x \in D_x \ \forall y \in D_y \ \varphi(x) \wedge p_A(x, y)
    \wedge R_{ATA}(x, y) \rightarrow p_A(x, r_{ATA}(x, y)) \\
  AFTA:\ &\forall x \in D_x \ \forall y \in D_y \ \varphi(x) \wedge p_A(x, y)
    \wedge R_{AFTA}(x, y) \rightarrow p_A(x, r_{AFTA}(x, y))
\end{align}

Условие корректоности оценочной функции:
\begin{align}
  CORR_A:\ &\forall x \in D_x \ \forall y \in D_y \ \varphi(x) \wedge p_A(x, y)
    \rightarrow u_A(x, y) \in W \\
\end{align}

Условие завершимости:
\begin{align}
\begin{split}
  ATA: \ &\forall x \in D_x \ \forall y \in D_y \ \varphi(x) \wedge p_A(x, y)
    \wedge R_{ATA}(x, y) \\
    &\quad \rightarrow u_A(x, r_{ATA}(x, y)) < u_A(x, y)
\end{split} \\
\begin{split}
  AFTA:\ &\forall x \in D_x \ \forall y \in D_y \ \varphi(x) \wedge p_A(x, y)
    \wedge R_{AFTA}(x, y) \\
    &\quad \rightarrow u_A(x, r_{AFTA}(x, y)) < u_A(x, y)
\end{split}
\end{align}

Возьмём:
\begin{align}
  & p_A(x, y) \equiv y_1 \geq x \wedge y_2 \geq 1 \wedge y_3 \geq x
    \wedge y_1 = xy_2 \wedge y_2 \leq y_3 \wedge y_3 \leq x^2 \\
  & W \equiv (\mathbb{N}_0,\mathbb{N}_0) \\
  & <: (a, b) < (c, d) \Leftrightarrow a < c \vee a = c \wedge b < d \\
  & u_A(x, y) = (x^2 - y_2, x^2 - y_3)
\end{align}

Получаем:
\begin{align}
  SA:    \ &\forall x \in D_x \ x > 1 \rightarrow x \geq x \wedge x \geq 1
    \wedge x \geq x \wedge x = 1x \wedge 1 \leq x \wedge x \leq x^2 \label{st3sa} \\
\begin{split}
  ATA:   \ &\forall x \in D_x \ \forall y \in D_y \ x > 1 \wedge y_1 \geq x
    \wedge y_2 \geq 1 \wedge y_3 \geq x \wedge y_1 = xy_2 \wedge y_2 \leq y_3 \\
    \wedge \ &y_3 \leq x^2 \wedge y_2 < y_3 \rightarrow y_1 + x \geq x
    \wedge y_2 + 1 \geq 1 \wedge y_3 \geq x \\
    \wedge \ &y_1 + x = x(y_2 + 1) \wedge (y_2 + 1) \leq y_3 \wedge y_3 \leq x^2
    \label{st3ata}
\end{split} \\
\begin{split}
  AFTA:  \ &\forall x \in D_x \ \forall y \in D_y \ x > 1 \wedge y_1 \geq x
    \wedge y_2 \geq 1 \wedge y_3 \geq x \wedge y_1 = xy_2 \wedge y_2 \leq y_3 \\
    \wedge \ &y_3 \leq x^2 \wedge y_2 \geq y_3 \wedge y_3 = x \\
    \rightarrow \ &y_1 \geq x \wedge y_2 \geq 1 \wedge y_1 \geq x
    \wedge y_1 = xy_2 \wedge y_2 \leq y_1 \wedge y_1 \leq x^2 \label{st3afta}
\end{split} \\
\begin{split}
  CORR_A:\ &\forall x \in D_x \ \forall y \in D_y \ x > 1 \wedge y_1 \geq x
    \wedge y_2 \geq 1 \wedge y_3 \leq y_3 \wedge y_3 \geq x \\
    \rightarrow \ &(x^2 - y_2, x^2 - y_3) \in (\mathbb{N}_0, \mathbb{N}_0)
    \label{st3corra}
\end{split} \\
\begin{split}
  ATA:   \ &\forall x \in D_x \ \forall y \in D_y \ x > 1 \wedge y_1 \geq x
    \wedge y_2 \geq 1 \wedge y_3 \geq x \wedge y_1 = xy_2 \wedge y_2 \leq y_3 \\
    \wedge \ &y_3 \leq x^2 \wedge y_2 < y_3
    \rightarrow (x^2 - y_2 - 1, x^2 - y_3) < (x^2 - y_2, x^2 - y_3) \label{st3atau}
\end{split} \\
\begin{split}
  AFTA:  \ &\forall x \in D_x \ \forall y \in D_y \ x > 1 \wedge y_1 \geq x
    \wedge y_2 \geq 1 \wedge y_3 \geq x \wedge y_1 = xy_2 \wedge y_2 \leq y_3 \\
    \wedge \ &y_3 \leq x^2 \wedge y_2 \geq y_3 \wedge y_3 = x
    \rightarrow (x^2 - y_2, x^2 - y_1) < (x^2 - y_2, x^2 - y_3) \label{st3aftau}
\end{split}
\end{align}

Пояснения к \cref{st3afta}:
\begin{align}
  & x > 1 \wedge y_1 = xy_2 \rightarrow y_2 \leq y_1 \\
  & x > 1 \wedge y_1 = xy_2 \wedge y_2 \leq y_3 \wedge y_3 = x
  \rightarrow y_1 \leq x^2
\end{align}

Пояснения к \cref{st3aftau}:
\begin{align}
  & x > 1 \wedge y_1 = xy_2 \wedge y_2 \geq y_3 \rightarrow y_1 > y_3
\end{align}


Из \cref{st3sa,st3ata,st3afta,st3corra,st3atau,st3aftau} вытекает
завершимость:
\begin{align}
  \{\varphi\}P\{\mathbb{T}\}
\end{align}

% Раскомментируйте, если нужно приложение
\appendix

\cleardoublepage \phantomsection
\section*{Приложение}\label{sec:appendix}
\addcontentsline{toc}{section}{Приложение}

\begin{align}\label{alt}
  \overline{AB} \equiv \overline{A} \vee \overline{B} \equiv A \rightarrow \overline{B}
\end{align}

\begin{align}\label{impl}
  A \rightarrow (B \rightarrow A) \equiv \overline{A} \vee \overline{B} \vee A
\end{align}

\begin{align}\label{any}
  \overline{A} \rightarrow (A \rightarrow B) \equiv A \vee \overline{A} \vee B
\end{align}

\begin{align}\label{hack}
  (\overline{A} \vee \overline{B}) \wedge B \rightarrow \overline{A} \equiv
  \overline{\overline{A} \vee \overline{B}} \vee \overline{A} \vee \overline{B}
\end{align}

\end{document}

